\documentclass{article}
\usepackage{url}

%%%%%%%%%%%% \linie %%%%%%%%%%%%%%%%%%%%%%%%%%%%%%%%
\newcommand{\linie}{\noindent {\rule{14cm}{1mm}}}
\pagenumbering{roman} \sloppy

\begin{document}

%%%%%%%%%%%%%%%%%%%%%%%%%%%%%%%%%%%%%%%%%%%%%%%%%%%%%%
%%%
%%%      The titlepage for ENTCS
%%%      (selfmade)
%%%%%%%%%%%%%%%%%%%%%%%%%%%%%%%%%%%%%%%%%%%%%%%%%%%%%%

\thispagestyle{empty} \vskip-2cm
\begin{sffamily} \noindent{\LARGE
 Electronic Notes in Theoretical Computer Science}\\
 \linie
 \vskip-.2cm
 \linie
 \vskip1cm
 \noindent \large Volume 82.1\\
 \vskip 1cm
 \noindent {\Large \sc  Coalgebraic Methods in Computer Science }\\

 \noindent {\Large \sc \bf CMCS'03} \ \ \vskip 1.2cm
 \noindent\large Warsaw, Poland\\

 \noindent April 5-6, 2003\\

 \vskip.3cm
 \linie
 \vskip7cm \noindent
 Guest Editor:\\

 \noindent {\sc H.~Peter Gumm} \vskip.5cm
 \linie
\end{sffamily}
\newpage

%%%%%%%%%%%%%%%%%%%%%%%%%%%%%%%%%%%%%%%%%%%%%%%%%%%%%%
%%%
%%%      An empty page
%%%
%%%%%%%%%%%%%%%%%%%%%%%%%%%%%%%%%%%%%%%%%%%%%%%%%%%%%%

\ \
\newpage


%%%%%%%%%%%%%%%%%%%%%%%%%%%%%%%%%%%%%%%%%%%%%%%%%%%%%%
%%%
%%%      Table of content for proceedings
%%%      using the proc-package
%%%
%%%%%%%%%%%%%%%%%%%%%%%%%%%%%%%%%%%%%%%%%%%%%%%%%%%%%%

%%%  An entry for the preface
\addtocontents{toc}{\contentsline {section}{ Preface}{\rm v}}

%%%%%%%%%%%%%%%%%%%%%%%%%%%%%%%%%%%%%%%%%%%%%%%%%%%%%%%%%%%%%%%
%%%
%%%      Process all entries from cmcs03db.tex
%%%
%%%%%%%%%%%%%%%%%%%%%%%%%%%%%%%%%%%%%%%%%%%%%%%%%%%%%%%%%%%%%%%

 \input proctoc         %% macros for generating each entry
 \input cmcs03db        %% database for the proceedings

%%%       Generate the table of contents :

 \tableofcontents

%%%%%%%%%%%%%%%%%%%%%%%%%%%%%%%%%%%%%%%%%%%%%%%%%%%%%%%%%%%%%%%%%
%%%%%
%%%%%   The Preface
%%%%%
%%%%%%%%%%%%%%%%%%%%%%%%%%%%%%%%%%%%%%%%%%%%%%%%%%%%%%%%%%%%%%%%
%
\newpage
\section*{Preface}
This volume contains the proceedings of the Sixth Workshop on
Coalgebraic Methods in Computer Science (CMCS'03). The workshop
was held in Warsaw, Poland on April 5 and 6, 2003, as a satellite
event to the European Joint Conference on Theory and Practice of
Software (ETAPS'03). The aim of the CMCS workshop series is to
bring together researchers with a common interest in the theory of
coalgebras and its applications. Previous workshops have been
organized in Lisbon (1998), Amsterdam (1999), Berlin (2000), Genoa
(2001), and Grenoble (2002).

Coalgebras have been found extremely useful for capturing
state-based dynamical systems, such as transition systems,
automata, process calculi, and class-based systems. The theory of
coalgebras has developed into a field of its own interest,
presenting a deep mathematical foundation, a growing domain of
applications and interactions with various other fields, such as
reactive and interactive systems theory, object oriented and
concurrent programming, formal system specification, modal logic,
dynamical systems, control systems, category theory, algebra,
analysis, etc..
%Much of this is witnessed by the papers in this volume.

 \vskip.5cm

\noindent The program committee of CMCS'03 consisted of
\begin{itemize}
 \item Ji\v r\'\i\ Ad\'amek, Technical University of
 Braunschweig, Germany;
 \item Corina C\^{\i}rstea, Oxford University, UK;
 \item H.~Peter Gumm, University of Marburg, Germany, (chair);
 \item Bart Jacobs, University of Nijmwegen, The Netherlands;
 \item Alexander Kurz, University of Leicester, UK;
  \item Marina Lenisa, University of Udine, Italy;
  \item Ugo Montanari, University of Pisa, Italy;
 \item Larry Moss, Indiana University, Bloomington, USA;
 \item Ataru T. Nakagawa, Software Research Associates, Tokyo
 Japan;
 \item Horst Reichel, Technical University Dresden, Germany;
 \item Grigore Ro\c{s}u, University of Illinois, Urbana, USA;
 \item Jan Rutten, CWI, Amsterdam, The Netherlands;
 \item James Worrell, Tulane University, New Orleans, USA.
\end{itemize}

\noindent The papers were refereed by the program committee and by
several outside referees, whose help is gratefully acknowledged.
The invited speakers at the conference were
\begin{itemize}
 \item Zolt\'{a}n \'{E}sik, University of Szeged, Hungary;
 \item Vaughan Pratt, Stanford University, USA.
\end{itemize}
Vaughan Pratt's lecture is included in the printed version of
these proceedings, Zolt\'{a}n \'{E}sik's lecture is expected to be
part of the final ENTCS volume.


These proceedings will be published as volume 82.1 in the series
Electronic Notes in Theoretical Computer Science (ENTCS). The
proceedings of the previous CMCS workshops appeared as ENTCS
Volumes 11, 19, 33, 44.1, and 65.1.

ENTCS is published electronically through the facilities of
Elsevier Science B.V. and under its auspices. The volumes in the
ENTCS series are available online at
 \url{http://www.elsevier.nl/locate/entcs}.
We are grateful to ENTCS for their continuing support, in
particular to Mike Mislove, Managing Editor of the ENTCS series.

For the sixth time, CMCS has been organized as satellite event to
ETAPS. We are very grateful to the ETAPS organizers, especially to
Damian Niwinski, for taking care of all the local organization and
for accommodating all our special requests. Thanks also to Mikolaj
Bojanczyk for managing the printing of these proceedings.\\

\noindent \emph{Marburg, February 28, 2003} \hfill \emph{H.~Peter
Gumm}

\end{document}
